\documentclass[10pt,a4paper]{amsart}
\usepackage[margin=19mm]{geometry}
\usepackage{amsmath,amssymb,mathtools}
\usepackage[scr=rsfs,cal=euler]{mathalfa}
\usepackage{xfrac}
\usepackage[unicode,hidelinks,bookmarks=false]{hyperref}
\newcommand{\ie}{i.e.\ }
\newcommand{\cf}{cf.\ }
\newcommand{\resp}{respectively\ }
\newcommand{\Subsection}{\subsection}
\providecommand{\nobreakdash}{\nobreak\hbox{-}}
\setlength{\emergencystretch}{2em}
\multlinegap0pt
\usepackage{amssymb}
\usepackage{enumerate}
\usepackage{xcolor}
\usepackage{mathtools}
\usepackage[shortlabels]{enumitem}
\newtheorem{cor}{Corollary}
\DeclareMathOperator{\Bl}{Bl}
\newcommand{\KdimS}{K_0(\Var_S^{\dim})}
\DeclareMathOperator{\Var}{Var}
\title{Involution on the Graded Grothendieck Ring of Varieties and $\mathbb{D}$-Singularities}
\author{Andrew Burke}
\date{Source: arXiv:2508.17587v2 (CC BY 4.0)}
\begin{document}
\maketitle

\noindent\emph{Source and licence.} Involution on the Graded Grothendieck Ring of Varieties and $\mathbb{D}$-Singularities. Version arXiv:2508.17587v2 (CC BY 4.0), distributed by arXiv under the Creative Commons Attribution 4.0 International licence, http://creativecommons.org/licenses/by/4.0/, as recorded in the arXiv licence metadata for this exact version and shown on the abstract page https://arxiv.org/abs/2508.17587v2. Full licence notice: \texttt{licenses/arxiv-2508.17587-CC-BY-4.0.txt}.\par\smallskip

\noindent\textit{Editorial scope.} Counted unit: the article's cor module (source0.tex lines 468--484). The article's complete original proof of it is delivered with it (source0.tex lines 486--511, one \texttt{proof} environment of 26 lines). No mathematics is written, repaired or re-derived: the statement and the proof are the article's own lines, in the article's own notation, and no retained line is substituted or re-flowed.  No further source-local context is needed: every cross-reference of every retained line resolves inside the two delivered spans.  Nothing was omitted from any retained span, so the package carries no omission record.  No retained line contains a citation, so the wrapper carries no bibliography and no bibliography entry.  No retained line contains a figure environment, an image-inclusion command or drawing code, so no image is delivered and the package declares no asset dependency.  Editorial formatting only, in addition: an AMS \texttt{amsart} preamble that restates the article's own declarations and macros used by the retained lines, namely the environments \texttt{cor} on separate counters, together with the standard packages those lines need.  The retained environments print in this document as Corollary 1 in order of appearance, whereas the article prints the same items with the numbering of its own full document.  The complete native source of the article as distributed in the arXiv e-print for this exact version is bundled unchanged as \texttt{sources/183796/source0.tex}.
\begin{cor}\label{bit module}
The graded Grothendieck group $\KdimS$ of $S$-varieties admits an alternative presentation as a $\mathbb{Z}[\tau, \mathbb{L}]$-module with:

\begin{itemize}
    \item generators in degree $d$ the isomorphism classes $[X]$ of regular, connected $S$-schemes $X$, which are proper, of finite type and of relative dimension $d$ over $S$, and
    \item relations
        \begin{align}\tag{1}
            [\Bl_Y X] = [X] + \tau \mathbb{L} \cdot (\tau^{k-1} + \ldots + \mathbb{L}^{k-1}) \cdot [Y]
        \end{align}
        for $\Bl_Y X$ the blow-up of a regular, connected $S$-scheme $X$, proper and of finite type over $S$ along a regular and connected subscheme $Y$ of codimension $k+1$, and with $E$ the exceptional divisor, and 
        \begin{align}\tag{2}
            [E] = (\tau^{k}+ \tau^{k-1} \mathbb{L} +  \ldots + \mathbb{L}^{k}) \cdot [Y]
        \end{align} 
        for $E \rightarrow Y$ a Zariski locally-trivial $\mathbb{P}^k$-fibration over a regular, connected $S$-scheme $Y$, proper and of finite type over $S$.
        
\end{itemize}
\end{cor}

\begin{proof}
Temporarily denote by $K_0(\Var_S^{\tau, \mathbb{L}})$ the $\mathbb{Z}[\tau, \mathbb{L}]$-module defined by the above presentation. We make use of the Bittner presentation on $\KdimS$. Consider $\Theta: \KdimS \rightarrow K_0(\Var_S^{\tau, \mathbb{L}})$ defined on generators by $[X]_d \mapsto [X] \cdot \tau^{d - \dim_S X}$. This is a well-defined group homomorphism, for
\begin{align*}
    \Theta([\Bl_Y X]_d) &- \Theta([E]_d) - \Theta([X]_d) + \Theta([Y]_d) \\
    &= \big( [\Bl_Y X] - \tau \cdot [E] -[X] + \tau^{k+1} \cdot [Y]  \big) \cdot \tau^{d-\dim_S X} \\
    &= \big( [\Bl_Y X] - \tau \cdot (\tau^{k}+ \ldots + \mathbb{L}^{k})[Y] -[X] + \tau^{k+1} \cdot [Y]  \big) \cdot \tau^{d-\dim_S X} \\
    &= \big( [\Bl_Y X] -[X] - \tau \mathbb{L} \cdot ( \tau^{k-1} + \ldots + \mathbb{L}^{k-1}) \cdot [Y]  \big) \cdot \tau^{d-\dim_S X} \\
    &= 0,
\end{align*}
if $\Bl_Y X \rightarrow X$ is the blow-up along a regular subscheme $Y$ and $d \geq \dim_S X$. It is moreover, a $\mathbb{Z}[\tau, \mathbb{L}]$-module homomorphism.

Define an inverse $\Theta^{-1}:  K_0(\Var_k^{\tau, \mathbb{L}}) \rightarrow \KdimS$ on generators by $[X] \mapsto [X]_{\dim_S X}$. To see that it is well-defined, suppose $E \rightarrow Y$ is a Zariski locally-trivial $\mathbb{P}^k$-fibration. Then
\begin{align*}
    \Theta^{-1}([E]) = [E]_{\dim_S E} = (\tau^k + \ldots + \mathbb{L}^k) \cdot 
 [Y]_{\dim_S Y} = \Phi^{-1}((\tau^k + \ldots + \mathbb{L}^k) \cdot [Y]).
\end{align*}

Moreover, if $\Bl_Y X \rightarrow X$ is the blow-up along a regular subscheme $Y$, we compute
\begin{align*}
    \Theta^{-1}&([\Bl_Y X]) - \Theta^{-1}([X]) - \Theta^{-1}(\tau \mathbb{L} \cdot (\tau^{k-1} + \ldots + \mathbb{L}^{k-1}) \cdot [Y]) \\
    =& [\Bl_Y X]_d - [X]_d - \big( \tau \cdot (\tau^k + \ldots + \mathbb{L}^k) - \tau^k \big) \cdot [Y]_{d-k-1} \\
    =& [\Bl_Y X]_d - [X]_d - [E]_{d} + [Y]_{d} \\
    =& 0.
\end{align*}
The result follows.
\end{proof}

\end{document}
